\section{Coordinates as a language}

The purpose of this chapter is not only to introduce several formulas. Its purpose is to build a language. Coordinates, points, vectors, lengths, distances, and dot products are the words of that language. Once we can read them correctly, we can use them to describe more complicated geometric objects.

The next chapters will study lines and planes. These objects will not appear from nowhere. They will be built from the ideas already introduced here.

\subsection*{Points give position}

A point tells us where we are. In the plane, a point has the form
\[
P=(x,y).
\]
In space, it has the form
\[
P=(x,y,z).
\]
A point may be fixed, such as
\[
A=(2,-1),
\]
or variable, such as
\[
P=(x,y).
\]
When the coordinates are variable, the point is allowed to move through a set of possible positions.

This distinction is essential for equations. An equation in coordinates usually does not describe one number. It describes all points whose coordinates satisfy a condition.

\begin{examplebox}
The condition
\[
y=3
\]
describes all points \((x,y)\) in the plane whose second coordinate is \(3\). The first coordinate is free. Thus the equation describes a horizontal line.

The point \((2,3)\) is one point on that line, but it is not the whole line.
\end{examplebox}

A common mistake is to confuse one solution with the entire geometric object. A single point may satisfy an equation, but the equation may describe a whole set of points.

\subsection*{Vectors give displacement and direction}

A vector tells us how to move. If
\[
v=(v_1,v_2),
\]
then \(v_1\) and \(v_2\) describe coordinate changes. The vector may be used as an actual displacement from one point to another, or it may be used only as a direction.

For example, the vector
\[
v=(2,1)
\]
can mean: move two units in the \(x\)-direction and one unit in the \(y\)-direction. But it can also mean: consider the direction determined by this displacement.

This distinction will matter for lines. A line can be described by one point on the line and one non-zero direction vector. The point anchors the line. The vector gives the direction in which the line extends.

\begin{remarkbox}
A point and a vector may have the same coordinates, but they play different roles. A point answers the question ``where?'' A vector answers the question ``how far and in which direction?''
\end{remarkbox}

\subsection*{Equations describe conditions}

An equation in analytic geometry is usually a condition on points. It says which points are allowed and which points are not.

For example, the equation
\[
x^2+y^2=1
\]
describes all points whose distance from the origin is \(1\). This is a circle of radius \(1\) centered at the origin.

The formula is algebraic, but the object is geometric. The equation does not merely ask us to solve for \(x\) or \(y\). It describes a set of points.

\begin{examplebox}
Consider the condition
\[
(x-2)^2+(y+1)^2=9.
\]
A point \((x,y)\) satisfies this condition exactly when its distance from \((2,-1)\) is \(3\). Therefore the equation describes a circle with center \((2,-1)\) and radius \(3\).
\end{examplebox}

This example shows the main power of analytic geometry. A geometric statement about distance becomes an algebraic equation. Conversely, an algebraic equation can be read as a geometric statement.

\subsection*{Fixed and variable quantities}

Many mistakes in analytic geometry come from not distinguishing fixed objects from variable ones. A point may be fixed. A direction vector may be fixed. But the point \((x,y)\) or \((x,y,z)\) often represents a variable point moving through a set.

When we write a line in the form
\[
P=A+t v,
\]
the point \(A\) is fixed, the vector \(v\) is fixed and non-zero, and the parameter \(t\) is variable. As \(t\) changes, the point \(P\) moves along the line.

This formula belongs to the next chapter, but the language is already available. A fixed point gives a starting position. A direction vector gives a possible movement. A parameter controls how far we move.

\begin{remarkbox}
Whenever a formula contains letters, ask which letters represent fixed data and which letters are allowed to vary. This question often reveals the meaning of the formula.
\end{remarkbox}

\subsection*{From this chapter to lines and planes}

The chapter has prepared three ideas that will be used repeatedly.

First, coordinates describe position. Second, vectors describe displacements and directions. Third, equations describe conditions on variable points.

Lines will be described by combining points and direction vectors. Planes will be described either by points and directions or by a point and a normal vector. The dot product will allow us to express perpendicularity, and therefore it will help us write equations of planes.

\begin{summarybox}
When working with coordinates and vectors, ask:
\begin{itemize}
\item Am I describing a point, a vector, or a set of points?
\item Which quantities are fixed, and which are variable?
\item Is a vector being used as a displacement, a direction, or a normal vector?
\item Can a geometric condition be translated into an algebraic equation?
\item After computing, can I interpret the result geometrically?
\end{itemize}
\end{summarybox}

Analytic geometry is not a collection of unrelated formulas. It is a way of translating between geometric objects and algebraic statements. The translation is useful only when we understand both languages.